% ========================================================
% TRABAJO EN GITHUB
% ========================================================
\section{Trabajo en GitHub}

El proyecto se aloja en el repositorio público \url{https://github.com/R0SEWT/concurrente}. El
Trabajo Parcial ocupa el directorio \texttt{tp/}, con su propio toolchain y una integración continua
en GitHub Actions que corre las pruebas del pipeline en cada pull request.

\subsection{Flujo Git Flow}

El trabajo sigue Git Flow (Tabla~\ref{tab:ramas}): ningún cambio entra directamente a
\texttt{develop} ni a \texttt{main}, sino por pull request desde una rama propia. La
Tabla~\ref{tab:prs} lista los pull requests e issues del Trabajo Parcial.

\begin{table}[H]
\caption{Ramas del flujo de trabajo}\label{tab:ramas}
\small
\begin{tabularx}{\textwidth}{@{}L{3.4cm}Y@{}}
\toprule
\textbf{Rama} & \textbf{Uso} \\
\midrule
\texttt{main} & Versión entregada; recibe una fusión desde \texttt{develop} por entregable \\
\texttt{develop} & Integración; todo cambio llega por pull request revisado \\
\texttt{feature/*}, \texttt{feat/*} & Una rama por unidad de trabajo, que sale de \texttt{develop} y
  vuelve a ella \\
\bottomrule
\end{tabularx}
\end{table}

\begin{table}[H]
\caption{Pull requests e issues del Trabajo Parcial}\label{tab:prs}
\small
\begin{tabularx}{\textwidth}{@{}lL{3.6cm}YL{2cm}@{}}
\toprule
\textbf{N.º} & \textbf{Autor} & \textbf{Contenido} & \textbf{Estado} \\
\midrule
\#3 & \nombreDayana & Propuesta de caso: SENAMHI y ODS 13 & Cerrada \\
\#4 & \nombreDayana & Propuesta de caso: NYC TLC y ODS 11 & Fusionada \\
\#5 & \nombreRody & Issue de investigación bibliográfica & Abierta \\
\#6 & \nombreRody & Limpieza del dataset & Fusionada \\
\#7 & \nombreDayana & Análisis exploratorio (notebook) & Fusionada \\
\#8 & \nombreRody & Bibliografía en BibTeX & Fusionada \\
\#9 & \nombreRody & CI: notebooks sin salidas & En revisión \\
\#16 & \nombreRody & Informe de PC1 en LaTeX & Fusionada \\
\bottomrule
\end{tabularx}
\end{table}

La versión entregada corresponde al tag \texttt{pc1} de la rama \texttt{main}.

\subsection{Historial de commits por integrante}

Las Tablas~\ref{tab:commits-autor} y~\ref{tab:commits} se generan desde el repositorio con
\texttt{./generar-historial.sh} y cuentan los commits que tocan el Trabajo Parcial en todas las
ramas publicadas.

% Generado por generar-historial.sh el 2026-09-13. No editar a mano.
\begin{table}[H]
\caption{Commits del Trabajo Parcial por integrante}\label{tab:commits-autor}
\small
\begin{tabular}{@{}lrrll@{}}
\toprule
\textbf{Integrante} & \textbf{Commits} & \textbf{Merges de PR} & \textbf{Primero} & \textbf{Último} \\
\midrule
Dayana Kety Gómez Rodriguez & 6 & 1 & 2026-09-10 & 2026-09-13 \\
Rody Sebastian Vilchez Marin & 17 & 0 & 2026-09-11 & 2026-09-13 \\
\bottomrule
\end{tabular}

{\footnotesize\textit{Nota.} Todas las ramas publicadas en \texttt{origin}. Los merges de PR se cuentan aparte porque no tocan archivos por sí mismos.\par}
\end{table}

{\small
\begin{xltabular}{\textwidth}{@{}L{2cm}L{2.9cm}L{1.5cm}Y@{}}
\caption{Historial de commits del Trabajo Parcial, del más reciente al más antiguo}\label{tab:commits}\\
\toprule
\textbf{Fecha} & \textbf{Autor} & \textbf{Commit} & \textbf{Mensaje} \\
\midrule
\endfirsthead
\toprule
\textbf{Fecha} & \textbf{Autor} & \textbf{Commit} & \textbf{Mensaje} \\
\midrule
\endhead
\bottomrule
\endfoot
2026-09-13 & Dayana Kety Gómez Rodriguez & \texttt{4ed330d} & docs: agregar conclusion individual Dayana \\
2026-09-13 & Rody Sebastian Vilchez Marin & \texttt{dfeb8fb} & Evitar la cita redundante en las leyes de Amdahl y Gustafson \\
2026-09-13 & Rody Sebastian Vilchez Marin & \texttt{a6b11b6} & Agregar la opinión sobre Voevodski y el análisis de Amdahl y Gustafson \\
2026-09-13 & Dayana Kety Gómez Rodriguez & \texttt{9a7a407} & docs: agregar opiniones críticas del tema 2 \\
2026-09-13 & Rody Sebastian Vilchez Marin & \texttt{8143e69} & Concretar el experimento de memoria en la opinión sobre Popcorn \\
2026-09-13 & Rody Sebastian Vilchez Marin & \texttt{cf415a7} & Agregar la opinión crítica sobre Popcorn (Bellavita et al., 2025) \\
2026-09-13 & Rody Sebastian Vilchez Marin & \texttt{df932c1} & Agregar la opinión crítica sobre Li et al. (2023) \\
2026-09-13 & Rody Sebastian Vilchez Marin & \texttt{54318a7} & Detallar la ficha de Li et al. (2023) con el texto completo \\
2026-09-13 & Rody Sebastian Vilchez Marin & \texttt{745cacc} & Abortar la tabla de commits si no hay refs de origin \\
2026-09-13 & Rody Sebastian Vilchez Marin & \texttt{f879cc9} & Actualizar la tabla de commits del informe \\
2026-09-13 & Rody Sebastian Vilchez Marin & \texttt{d640e77} & Redactar el informe de PC1 en LaTeX \\
2026-09-12 & Rody Sebastian Vilchez Marin & \texttt{e1c2733} & Dejar los 9 papers de la cuota en acceso abierto \\
2026-09-12 & Rody Sebastian Vilchez Marin & \texttt{daab286} & Incorporar los papers de la carpeta de Drive del equipo \\
2026-09-12 & Rody Sebastian Vilchez Marin & \texttt{37e66a6} & Verificar en CI que los notebooks no lleguen con salidas \\
2026-09-12 & Rody Sebastian Vilchez Marin & \texttt{5a7472f} & Corregir el año de la entrada he2022 (CCPE) \\
2026-09-11 & Rody Sebastian Vilchez Marin & \texttt{650c7b0} & Merge branch 'develop' into feature/bibliografia-pc1 \\
2026-09-11 & Rody Sebastian Vilchez Marin & \texttt{7675cfb} & Agregar la bibliografía de PC1 en BibTeX \\
2026-09-11 & Dayana Kety Gómez Rodriguez & \texttt{fbec506} & feat: implement structural updates and optimizations across multiple modules \\
2026-09-11 & Rody Sebastian Vilchez Marin & \texttt{a6fead9} & Silenciar el falso positivo de opengrep en la auditoría DuckDB \\
2026-09-11 & Rody Sebastian Vilchez Marin & \texttt{7254176} & Tolerar silver vacío y parametrizar el SQL de la auditoría \\
2026-09-11 & Rody Sebastian Vilchez Marin & \texttt{9617a56} & Limpiar dataset NYC TLC yellow 2024-01 (bronze → silver → gold) \\
2026-09-10 & Dayana Kety Gómez Rodriguez & \texttt{e2a4a6f} & docs: agregar enlace a dataset de taxis amarillos de Azure \\
2026-09-10 & Dayana Kety Gómez Rodriguez & \texttt{8d7a2d8} & docs: agregar propuesta de caso de uso NYC TLC ODS 11 \\
2026-09-10 & Dayana Kety Gómez Rodriguez & \texttt{3c55c07} & docs: agregar propuesta de caso de uso SENAMHI ODS 13 \\
\end{xltabular}
}


El tercer integrante no registra commits en el repositorio.
